% ============================================================================
% WO-VSEPR_SIM-63A-ARCHIVE
% Day 63 Archive: Hierarchical Sampling and Convergence Testing
% ============================================================================
\documentclass[11pt, a4paper, oneside]{article}

\usepackage[top=1.0in, bottom=1.0in, left=1.15in, right=1.15in, headheight=14pt]{geometry}
\usepackage[T1]{fontenc}
\usepackage{lmodern}
\usepackage{microtype}
\usepackage{amsmath, amssymb, mathtools}
\usepackage{bm}
\usepackage{booktabs}
\usepackage{array}
\usepackage{tabularx}
\usepackage{longtable}
\usepackage{enumitem}
\setlist{nosep, leftmargin=1.5em}
\usepackage{fancyhdr}
\usepackage{titlesec}
\usepackage[colorlinks=true, linkcolor=black, citecolor=black, urlcolor=blue]{hyperref}

\pagestyle{fancy}
\fancyhf{}
\fancyhead[L]{\small\textit{VSEPR-SIM Work Order Archive}}
\fancyhead[R]{\small\textit{WO-VSEPR\_SIM-63A}}
\fancyfoot[C]{\thepage}
\renewcommand{\headrulewidth}{0.4pt}

\titleformat{\section}{\Large\bfseries}{\thesection}{1em}{}
\titleformat{\subsection}{\large\bfseries}{\thesubsection}{1em}{}

\newcommand{\R}{\mathbb{R}}
\DeclareMathOperator{\CV}{CV}

\begin{document}

\begin{center}
  {\LARGE\bfseries WO-VSEPR\_SIM-63A-ARCHIVE}\\[0.4em]
  {\Large\bfseries Hierarchical Sampling and Convergence Testing for High-$N$ Bead Simulation}\\[1.0em]
  {\normalsize VSEPR-SIM Development Day 63}\\[0.3em]
  {\small Archive Draft}
\end{center}

\vspace{0.5em}
\hrule
\vspace{1.2em}

\begin{abstract}
This work order archive defines a practical hierarchy for high-$N$ bead-resolved simulation in VSEPR-SIM. Direct all-pairs interaction evaluation scales as $O(N^2)$ and becomes infeasible for million-bead regimes. The document uses octree decomposition as the computational baseline and introduces a dynamic tree sampling method in which local subdivision depth and measurement density adapt to bead concentration, variance, defect concentration, and convergence instability. The purpose is constrained: not to claim atomistic electronic truth from encoded beads, but to test whether bead-level interaction rules converge toward stable macro-scale measurements under increasing sample count and representative-volume scale. Limitations and validation gates are made explicit to prevent overclaiming.
\end{abstract}

\tableofcontents
\newpage

\section{Introduction}

Pairwise interaction models scale quadratically:
\begin{equation}
\text{cost}_{\text{all-pairs}} = O(N^2)
\end{equation}
For $N=10^6$ beads, direct evaluation implies approximately
\begin{equation}
N(N-1)/2 \approx 5\times 10^{11}
\end{equation}
interaction checks per step (order of $10^{12}$ operations when force kernels and bookkeeping are included).

This is the core bottleneck in large-bead simulation.

\subsection{Why high-\texorpdfstring{$N$}{N} is still valuable}

Low-$N$ runs reveal local structure and implementation behavior. High-$N$ runs are needed for:
\begin{itemize}
  \item distribution-level measurements,
  \item finite-size convergence testing,
  \item representative-volume behavior,
  \item macro-property stability under repeated seeds.
\end{itemize}

\textbf{Key statement:} high-$N$ improves statistical convergence, not microscopic truth by itself.

\section{Established Octree Method}

In 3D, each cube is subdivided into $2\times 2\times 2 = 8$ child cubes.
This gives the canonical octree hierarchy.

The force decomposition is:
\begin{equation}
\mathbf{F}_i = \mathbf{F}_{i,\mathrm{near}} + \mathbf{F}_{i,\mathrm{far}}
\end{equation}
Near-field terms are evaluated explicitly; far-field terms are approximated via aggregate node descriptors.

The Barnes--Hut acceptance condition is:
\begin{equation}
\frac{s}{d} < \theta
\end{equation}
where:
\begin{center}
\begin{tabular}{ll}
\toprule
Symbol & Meaning \\
\midrule
$s$ & node size \\
$d$ & distance from target bead to node center \\
$\theta$ & approximation tolerance \\
\bottomrule
\end{tabular}
\end{center}

This relation provides the computational backbone: node aggregation is accepted only when geometric separation supports bounded approximation error.

\section{VSEPR-SIM High-\texorpdfstring{$N$}{N} Bead Interpretation}

To prevent overclaiming, the simulated units are defined as \emph{bead-resolved state carriers}, not guaranteed fully resolved quantum atoms.

A bead state may encode:
\begin{equation}
\{m,\mathbf{r},q,\epsilon,\sigma,T,\mathbf{x},\mathbf{v},\text{state},\text{material\_tag}\}
\end{equation}

High-$N$ targets improvement in stable measured properties such as:
\begin{equation}
\rho(\mathbf{x}),\; \phi,\; D_{\mathrm{eff}},\; \text{coordination},\; \text{defect fraction},\; \text{local variance},\; \text{representative-volume stability}
\end{equation}

These outputs are interpreted as convergence behavior of the encoded model, not direct proof of chemical or electronic correctness.

\section{Dynamic Tree Method (Day 63 Contribution)}

The standard octree is geometry-driven and mostly static in refinement policy.
The dynamic tree method extends this by making local refinement and sampling adaptive to measurement quality.

\subsection{Adaptive refinement criteria}

Possible split triggers for node $\nu$:
\begin{align}
 n_{\nu} &> n_{\max} \\
 \sigma_{\rho}(\nu) &> \tau_{\rho} \\
 \sigma_{P}(\nu) &> \tau_{P} \\
 \lvert \nabla P(\nu) \rvert &> \tau_{\nabla} \\
 E_{\mathrm{local}}(\nu) &> \tau_{E}
\end{align}
where $P$ is a sampled property (porosity, density, coordination, defect fraction, or other tracked field).

\subsection{Dynamic behavior rules}

\begin{center}
\begin{tabularx}{\linewidth}{lX}
\toprule
Condition & Action \\
\midrule
too many beads in node & subdivide \\
low variation & keep coarse \\
high defect variance & refine \\
high boundary/edge error & resample or reject \\
stable measurement & stop refining \\
unstable measurement & increase samples \\
\bottomrule
\end{tabularx}
\end{center}

This reframes the method as \textbf{adaptive hierarchical sampling for bead-resolved material inference}.

\section{Why Not Purely Random Cube Sampling}

Random cube probes are useful for analysis,
\begin{equation}
C_k = \mathrm{cube}(\mathbf{x}_k, s_k)
\end{equation}
but unconstrained random size creates scale bias because $V=s^3$ and large cubes dominate aggregate statistics.

Recommended design:
\begin{equation}
\text{fixed scale bands} + \text{random centers} + \text{convergence reporting}
\end{equation}

Example scale set:
\begin{equation}
s \in \{5,10,25,50,100,250,500\}
\end{equation}

For each scale $s$, report:
\begin{equation}
\mu_P(s),\; \sigma_P(s),\; \CV_P(s)
\end{equation}
with coefficient of variation
\begin{equation}
\CV_P(s) = \frac{\sigma_P(s)}{\mu_P(s)}
\end{equation}

This measures stability directly rather than smoothing away heterogeneity.

\section{Limitations in Goodness Fitting}

The method cannot prove that encoded beads are physically correct. It can test whether the encoded bead model converges to stable measured properties under increased sample size and scale.

\begin{center}
\begin{tabularx}{\linewidth}{lX}
\toprule
Limitation & Meaning \\
\midrule
bead encoding quality & poor bead rules can converge to stable but unphysical output \\
interaction law choice & force law determines what is physically represented \\
sampling scale & too small over-amplifies local noise \\
large-volume dominance & large samples can hide local critical variation \\
smoothing risk & smoothing can erase true heterogeneity \\
parameter explosion & too many fit knobs can become non-identifiable \\
validation shortage & without empirical anchors, no truth claim is justified \\
\bottomrule
\end{tabularx}
\end{center}

\textbf{Core sentence:} high-$N$ improves convergence of the model, not correctness of the model.

\section{Proposed Convergence Gates}

Use explicit acceptance gates instead of high-dimensional ad hoc goodness functions.
A property $P$ is scale-stable if:
\begin{equation}
\CV_P(s) < \tau
\end{equation}
and
\begin{equation}
\left|\frac{\mu_P(s_{k+1})-\mu_P(s_k)}{\mu_P(s_k)}\right| < \epsilon
\end{equation}
across multiple random seeds.

Example defaults:
\begin{equation}
\tau = 0.05,\qquad \epsilon = 0.03,\qquad \text{seed\_repeat\_count}=3
\end{equation}

This yields a conservative and defensible validation layer.

\subsection{Gate decision table}

\begin{center}
\begin{tabularx}{\linewidth}{lX}
\toprule
Gate result & Interpretation \\
\midrule
$\CV_P(s) < \tau$ and drift $<\epsilon$ across seeds & Scale-stable for property $P$ at tested bands \\
$\CV_P(s) < \tau$ but drift $\ge \epsilon$ & Local variance controlled, but inter-scale consistency unresolved \\
$\CV_P(s) \ge \tau$ and drift $<\epsilon$ & Mean appears stable while fluctuations remain too high \\
$\CV_P(s) \ge \tau$ and drift $\ge \epsilon$ & Not converged; refine tree, increase samples, or revisit encoding \\
\bottomrule
\end{tabularx}
\end{center}

\subsection{Seed-robust pass policy}

For each property $P$, accept only if at least $2/3$ seed runs pass both gates
on the two largest scale bands. This avoids false positives caused by
single-seed luck in heterogeneous regions.

\section{Day 63 Framing}

Day 63 scope should prioritize hierarchical sampling and convergence testing, not new physics extensions.
The octree path supplies scalable structure; the dynamic tree layer supplies adaptive measurement discipline.

\subsection{Update-loop doctrine (Day 63+)}

Instead of a permanently one-way pipeline
\begin{equation}
\text{property prediction} \rightarrow \text{bead dynamics}
\end{equation}
the Day 63+ direction is a recursive coupling:
\begin{equation}
S_t \rightarrow P_t \rightarrow D_t \rightarrow S_{t+\Delta t}
\end{equation}
with
\begin{center}
\begin{tabular}{ll}
\toprule
Symbol & Meaning \\
\midrule
$S_t$ & bead state at time $t$ (positions, velocities, identities, local structure) \\
$P_t$ & inferred property field at time $t$ (density, porosity, coordination, transport proxy) \\
$D_t$ & bead-dynamics operator using inferred properties \\
$S_{t+\Delta t}$ & updated bead state \\
\bottomrule
\end{tabular}
\end{center}

Formally:
\begin{align}
P_t &= I(S_t) \\
S_{t+\Delta t} &= D(S_t, P_t) \\
P_{t+\Delta t} &= I(S_{t+\Delta t})
\end{align}

This is not circular reasoning as validation; it is recursive state--property coupling as dynamics.

\subsection{Paper-safe wording and non-claim boundary}

In this framework, circularity is not used as validation; it is used as coupling.
VSEPR-SIM infers local material descriptors from the bead-resolved state,
then bead dynamics uses those descriptors to update the system.
The updated state is then reanalyzed, producing a recursive state-property loop.
This loop is valuable because many material processes are path-dependent:
structure determines properties, while those properties govern the next structural transition.

\medskip
\noindent\textbf{Claim hygiene:}
\begin{itemize}
  \item Bad claim: the model proves its own properties because the dynamics reproduce them.
  \item Good claim: the model evolves self-consistent property fields under encoded bead rules.
  \item Better claim: the framework tests whether recursive bead-property coupling converges,
	destabilizes, or produces path-dependent material behavior.
\end{itemize}

\noindent\textbf{Day 63+ doctrine:}
\begin{quote}
Circularity is forbidden as validation, but essential as dynamics.\\
The loop is not proof. The loop is the model.
\end{quote}

Deliverables should therefore emphasize:
\begin{itemize}
  \item high-$N$ spatial decomposition and traversal behavior,
  \item scale-band convergence reporting,
  \item multi-seed stability checks,
  \item explicit non-claim boundaries.
\end{itemize}

\subsection{Intentional Day 63 bead-model simplicity}

For this paper/method, the bead model remains intentionally simple:
\begin{enumerate}
  \item spherical beads,
  \item assigned mass/radius/material tag,
  \item optional internal state,
  \item local interaction law,
  \item sampled property inference,
  \item convergence reporting,
  \item cleanup and reproducibility pass.
\end{enumerate}

\section{Implementation Notes for VSEPR-SIM}

A minimal integration path:
\begin{enumerate}
  \item Add dynamic refinement thresholds to runtime configuration.
  \item Record per-node variance and sample stability counters.
  \item Export scale-band summaries as structured report tables.
  \item Gate pass/fail by $(\CV,\epsilon)$ constraints across seeds.
\end{enumerate}

Suggested report columns per property and scale:
\begin{equation}
\{s,\;N_{\text{samples}},\;\mu_P,\;\sigma_P,\;\CV_P,\;\Delta\mu_{k\to k+1},\;\text{pass/fail}\}
\end{equation}

\subsection{Runtime pseudocode}

\begin{verbatim}
for each seed in seeds:
	build root node from domain bounds
	while refinement_queue not empty:
		node = pop(refinement_queue)
		sample_property(node)
		if split_trigger(node):
			subdivide(node)
			push(children)
		else:
			commit(node)

for each property P and each scale band s:
	compute mu_P(s), sigma_P(s), CV_P(s)
	compute drift = |mu_P(s[k+1]) - mu_P(s[k])| / mu_P(s[k])
	evaluate convergence gates

aggregate seed pass rates
emit archive report
\end{verbatim}

\subsection{Minimal archive artifact set}

Recommended outputs for reproducibility:
\begin{itemize}
  \item Node statistics dump: occupancy, split depth, trigger reason.
  \item Scale-band summary table for each tracked property.
  \item Seed-level pass/fail ledger with gate values.
  \item Run manifest: simulation hash, threshold set, and timestamp.
\end{itemize}

\section{Particle Roles and Transient Carrier Dynamics}

\subsection{Role decomposition}

The bead model is extended with multiple particle classes using different
update rules, force sensitivities, and event behavior.

\begin{center}
\begin{tabularx}{\linewidth}{l l X}
\toprule
Class & Role & Primary behavior \\
\midrule
Beads & heavy state carriers & define material structure and long-range field landscape \\
Transient charged carriers & light/mobile carriers & mediate charge redistribution, local EM effects, capture/emission \\
Transient neutral carriers & uncharged mobile carriers & ballistic transport with stochastic collision/absorption/scattering \\
\bottomrule
\end{tabularx}
\end{center}

This gives a clean three-level architecture: heavy structural state,
mobile charged response, and mobile neutral collision response.

\subsection{Beads as heavy structural carriers}

Each bead $i$ stores
\begin{equation}
\mathbf{B}_i=(\mathbf{x}_i,\mathbf{v}_i,m_i,r_i,q_i,E_i,s_i)
\end{equation}
with position, velocity, mass, interaction radius, charge,
internal energy state, and material/state label.

Beads provide octree aggregates such as
$M_{\mathrm{node}}$, $Q_{\mathrm{node}}$, center-of-mass,
and center-of-charge descriptors.

\subsection{Transient charged carriers}

Define charged carrier state
\begin{equation}
\mathbf{C}_k=(\mathbf{x}_k,\mathbf{v}_k,m_k,q_k,E_k)
\end{equation}
with light mass, high mobility, and subcycled time updates.

Electrostatic response:
\begin{equation}
\mathbf{F}^{\mathrm{EM}}_k=q_k\,\mathbf{E}(\mathbf{x}_k)
\end{equation}
where $\mathbf{E}$ may include bead-octree fields,
nearby charged-carrier corrections, external fields,
and boundary contributions.

Capture on bead $i$ applies deltas:
\begin{equation}
q_i^{\mathrm{new}}=q_i+q_k,
\qquad
m_i^{\mathrm{new}}=m_i+m_k
\end{equation}
Mass contribution from electron-like carriers may be neglected if
the model is configured for charge-dominant updates.

\subsection{Transient neutral carriers}

Define neutral carrier state
\begin{equation}
\mathbf{N}_\ell=(\mathbf{x}_\ell,\mathbf{v}_\ell,m_\ell,E_\ell),
\qquad q_\ell=0
\end{equation}
so electrostatic forcing is disabled:
\begin{equation}
\mathbf{F}^{\mathrm{EM}}_\ell=0
\end{equation}

Neutral motion is modeled as ballistic transport plus stochastic
collision/scattering/absorption events with beads.
At this scale, default behavior is \textbf{not} gravity-tree driven.

\subsection{Multi-time-stepping and tree strategy}

Use a slow bead step and fast transient substeps:
\begin{equation}
\Delta t_B=1.0\,\mathrm{fs},
\qquad
\Delta t_C=0.01\,\mathrm{fs},
\qquad
K=\frac{\Delta t_B}{\Delta t_C}=100
\end{equation}

Operational policy:
\begin{itemize}
  \item rebuild bead octree each bead step,
  \item use local overlay structures for transient substeps
	(cell list / bucket grid / local map),
  \item store event deltas in pending logs,
  \item apply deltas at bead step boundary (baseline) or
	propagate deltas up the node path (higher temporal accuracy).
\end{itemize}

Hybrid trigger:
\begin{equation}
\text{if capture count small }\Rightarrow\text{ propagate deltas;
else trigger early rebuild.}
\end{equation}

\subsection{Force/effect method by pair type}

\begin{center}
\begin{tabularx}{\linewidth}{l l X}
\toprule
Interaction & Method & Notes \\
\midrule
Bead $\leftrightarrow$ bead bonded & fixed bond list & springs/polymer/structural skeleton \\
Bead $\leftrightarrow$ bead short-range & neighbor list & LJ/contact/adhesion/damping \\
Bead $\leftrightarrow$ bead long-range & octree & electrostatic or configured long-range field \\
Charged carrier $\leftrightarrow$ bead & bead tree + local correction & distant bead charge via tree \\
Charged carrier $\leftrightarrow$ charged carrier & local neighbor list & avoid default $O(N_C^2)$ unless tiny $N_C$ \\
Neutral carrier $\leftrightarrow$ bead & stochastic collision model & ballistic + event model \\
Neutral carrier $\leftrightarrow$ neutral carrier & usually omitted & include only for dense neutral flux \\
Carrier $\leftrightarrow$ boundary & event policy & escape/reflection/absorption by setup \\
\bottomrule
\end{tabularx}
\end{center}

\subsection{Capture, emission, and reaction events}

For carrier $p$ entering bead interaction region
$\|\mathbf{x}_p-\mathbf{x}_i\|<R_i^{\mathrm{int}}$,
evaluate event probability by either rate model
\begin{equation}
P_{\mathrm{capture}}=1-\exp(-\lambda_i\Delta t)
\end{equation}
or path-length model
\begin{equation}
P_{\mathrm{capture}}=1-\exp(-\Sigma_i\Delta\ell)
\end{equation}
then draw $u\sim U(0,1)$ and accept event if $u<P_{\mathrm{capture}}$.

Charged-capture update:
\begin{equation}
q_i^{\mathrm{new}}=q_i+q_p,
\quad
m_i^{\mathrm{new}}=m_i+m_p,
\quad
E_i^{\mathrm{new}}=E_i+E_p+\Delta E_{\mathrm{bind}}
\end{equation}

Neutral-capture update:
\begin{equation}
q_i^{\mathrm{new}}=q_i,
\quad
m_i^{\mathrm{new}}=m_i+m_p,
\quad
E_i^{\mathrm{new}}=E_i+E_p+\Delta E_{\mathrm{reaction}}
\end{equation}
and optional state transform
\begin{equation}
s_i^{\mathrm{new}}=\mathcal{R}(s_i,p)
\end{equation}

Emission event model:
\begin{equation}
P_{\mathrm{emit}}=1-\exp\!\bigl(-\lambda_{\mathrm{emit}}(E_i,s_i)\,\Delta t\bigr)
\end{equation}
with conservation updates
\begin{equation}
q_i^{\mathrm{new}}=q_i-q_p,
\quad
m_i^{\mathrm{new}}=m_i-m_p,
\quad
E_i^{\mathrm{new}}=E_i-E_p-E_{\mathrm{cost}}
\end{equation}
and emitted velocity initialization
\begin{equation}
\mathbf{v}_p=v_0\,\hat{\mathbf{n}}
\end{equation}
where $\hat{\mathbf{n}}$ is isotropic or field-biased.

\subsection{Simulation loop and logging doctrine}

Reference loop:
\begin{verbatim}
for each bead timestep:
	rebuild bead octree
	rebuild bead neighbor list if needed

	compute bead-bead bonded/short/long-range forces

	for each transient substep:
		update charged carrier forces from bead tree
		update charged local interactions
		update neutral ballistic motion

		check capture, emission, collision, boundary events
		store pending bead mass/charge/energy deltas
		optionally propagate charge deltas up tree

	apply pending bead updates
	integrate bead positions and velocities
	write diagnostics and event log
\end{verbatim}

Event records should be separated from state trajectory outputs.
State files store evolution; event logs store transformations.

\begin{verbatim}
{
  "time": 125.04,
  "event": "charged_capture",
  "bead_id": 48291,
  "particle_id": 903,
  "delta_q": -1.0,
  "delta_m": 0.00054858,
  "delta_E": 2.31,
  "position": [4.2, 8.1, 0.5]
}
\end{verbatim}

\subsection{Naming and claim safety}

Section naming should use
\textbf{Transient Charged and Neutral Carrier Particles},
with electron-like or neutron-like behavior treated as special cases.

This preserves physical caution while retaining the event-driven
architecture required for coupled bead-carrier dynamics.

\subsection{Signed particle-type assignments (negative codes)}

For implementation convenience, transient and special particle classes may use
\textbf{negative type assignments}. These are \emph{classification codes},
not physical sign assignments for mass or energy.

\begin{center}
\begin{tabular}{lll}
\toprule
Type code & Particle class & Notes \\
\midrule
$+1, +2, \ldots$ & Beads (material-resolved) & Primary heavy structural carriers \\
$-1$ & Charged transient (electron-like) & Fast mobile charge carrier \\
$-2$ & Charged transient (ion-like) & Optional positive/negative by $q$ field \\
$-3$ & Neutral transient (neutron-like) & Ballistic + stochastic collision model \\
$-4$ & Neutral transient (gamma-like packet) & Event/transport packet; no rest-mass update by default \\
$-5$ & Heavy transient (alpha-like packet) & Composite carrier/event token \\
$-6$ and below & Reserved special channels & Decay products, diagnostics, boundary-injected particles \\
\bottomrule
\end{tabular}
\end{center}

Recommended rule set:
\begin{itemize}
  \item Type sign determines class family (positive = bead, negative = transient/special).
  \item Physical properties remain explicit fields: $(m, q, E, s)$.
  \item Neutrality is encoded by $q=0$, not by type sign alone.
  \item Gamma-like behavior should be configured as an event/packet model,
	not interpreted as a bead mass contribution unless explicitly enabled.
\end{itemize}

This convention simplifies routing and pair-dispatch logic in kernels:
\begin{equation}
\texttt{if type > 0 => bead path; if type < 0 => transient/special path}
\end{equation}

while preserving physically meaningful state updates through explicit
per-particle fields.

\section{Conclusion}

WO-VSEPR\_SIM-63A is positioned as a methodological archive for scalable inference in bead-resolved simulation.
The contribution is practical: combine established octree acceleration with adaptive sampling and convergence gates to decide whether macro-properties are stable under scale and seed variation.

The document intentionally separates:\
(1) computational convergence claims from\
(2) physical truth claims.

That separation is required for technically defensible high-$N$ reporting.

\vfill
\begin{center}
\rule{0.45\linewidth}{0.4pt}\\[0.5em]
\small WO-VSEPR\_SIM-63A-ARCHIVE\\
\small Hierarchical sampling and convergence testing first.
\end{center}

\end{document}
